% GENERATED FILE. Edit canonical lessons, then run npm run generate:books.
% canonical-source-hash: fnv1a64:70353d3f53ffdff4
% canonical-lessons: FR-C56-triste, FR-C56-content, FR-C56-fille, FR-C56-grand-pere, FR-C56-grand-mere

\chapter{Sad, Pleased, Daughter, Grandfather, Grandmother}
\label{ch:fr-sad-pleased-daughter-grandfather-grandmother}
\hlchaptermodality{56}

\begin{chapteropening}
\textbf{By the end of this chapter:} \emph{I can name sad, pleased, a daughter, a grandfather, a grandmother.}

Complete the last lesson of chapter 56: i can name sad, pleased, a daughter, a grandfather, a grandmother.
\end{chapteropening}

\section[triste]{triste — sad}

\label{lesson:FR-C56-triste}



\begin{quote}
Before the new one: say the French for sadness, then the French for happy.
\end{quote}

\subsection*{You'll want to know: triste}
\textbf{triste} — \textquotedblleft{}sad\textquotedblright{}. The adjective beside \textbf{la tristesse}.

\begin{grammarlens}[title={What you've built}]
One more word for this chapter.
\end{grammarlens}

\subsection*{Your turn}
\begin{itemize}
\raggedright
  \item \emph{Say it:} \emph{triste}
  \item \emph{Say it:} \emph{triste}, once more
  \item \emph{Say it:} say \emph{triste}
  \item \emph{Recall:} say the French for sadness, then the French for happy, then say \emph{triste} again
\end{itemize}

\begin{tcolorbox}[breakable,colback=teal!4,colframe=teal!35!black,title={Before you move on}]
What does triste mean? (\textquotedblleft{}Sad\textquotedblright{}.) Say it once more.
\end{tcolorbox}
\section[content]{content — pleased, glad}

\label{lesson:FR-C56-content}



\begin{quote}
Before the new one: say the French for happy, then the French for sad.
\end{quote}

\subsection*{You'll want to know: content}
\textbf{content} — \textquotedblleft{}pleased, glad\textquotedblright{}.

\begin{grammarlens}[title={What you've built}]
One more word for this chapter.
\end{grammarlens}

\subsection*{Your turn}
\begin{itemize}
\raggedright
  \item \emph{Say it:} \emph{content}
  \item \emph{Say it:} \emph{content}, once more
  \item \emph{Say it:} say \emph{content}
  \item \emph{Recall:} say the French for happy, then the French for sad, then say \emph{content} again
\end{itemize}

\begin{tcolorbox}[breakable,colback=teal!4,colframe=teal!35!black,title={Before you move on}]
What does content mean? (\textquotedblleft{}Pleased, glad\textquotedblright{}.) Say it once more.
\end{tcolorbox}
\section[la fille]{la fille — a daughter; a girl}

\label{lesson:FR-C56-fille}



\begin{quote}
Before the new one: say the French for sad, then the French for pleased.
\end{quote}

\subsection*{You'll want to know: la fille}
\textbf{la fille} — \textquotedblleft{}a daughter; a girl\textquotedblright{}.

\begin{grammarlens}[title={What you've built}]
One more word for this chapter.
\end{grammarlens}

\subsection*{Your turn}
\begin{itemize}
\raggedright
  \item \emph{Say it:} \emph{la fille}
  \item \emph{Say it:} \emph{la fille}, once more
  \item \emph{Say it:} say \emph{la fille}
  \item \emph{Recall:} say the French for sad, then the French for pleased, then say \emph{la fille} again
\end{itemize}

\begin{tcolorbox}[breakable,colback=teal!4,colframe=teal!35!black,title={Before you move on}]
What does la fille mean? (\textquotedblleft{}A daughter; a girl\textquotedblright{}.) Say it once more.
\end{tcolorbox}
\section[le grand-père]{le grand-père — a grandfather}

\label{lesson:FR-C56-grand-pere}



\begin{quote}
Before the new one: say the French for pleased, then the French for a daughter; a girl.
\end{quote}

\subsection*{You'll want to know: le grand-père}
\textbf{le grand-père} — \textquotedblleft{}a grandfather\textquotedblright{}.

\begin{grammarlens}[title={What you've built}]
One more word for this chapter.
\end{grammarlens}

\subsection*{Your turn}
\begin{itemize}
\raggedright
  \item \emph{Say it:} \emph{le grand-père}
  \item \emph{Say it:} \emph{le grand-père}, once more
  \item \emph{Say it:} say \emph{le grand-père}
  \item \emph{Recall:} say the French for pleased, then the French for a daughter; a girl, then say \emph{le grand-père} again
\end{itemize}

\begin{tcolorbox}[breakable,colback=teal!4,colframe=teal!35!black,title={Before you move on}]
What does le grand-père mean? (\textquotedblleft{}A grandfather\textquotedblright{}.) Say it once more.
\end{tcolorbox}
\section[la grand-mère]{la grand-mère — a grandmother}

\label{lesson:FR-C56-grand-mere}



\begin{quote}
Before the new one: say the French for a daughter; a girl, then the French for a grandfather.
\end{quote}

\subsection*{You'll want to know: la grand-mère}
\textbf{la grand-mère} — \textquotedblleft{}a grandmother\textquotedblright{}.

\begin{grammarlens}[title={What you've built}]
One more word for this chapter.
\end{grammarlens}

\subsection*{Your turn}
\begin{itemize}
\raggedright
  \item \emph{Say it:} \emph{la grand-mère}
  \item \emph{Say it:} \emph{la grand-mère}, once more
  \item \emph{Say it:} say \emph{la grand-mère}
  \item \emph{Recall:} say the French for a daughter; a girl, then the French for a grandfather, then say \emph{la grand-mère} again
\end{itemize}

\begin{tcolorbox}[breakable,colback=teal!4,colframe=teal!35!black,title={Before you move on}]
What does la grand-mère mean? (\textquotedblleft{}A grandmother\textquotedblright{}.) Say it once more.
\end{tcolorbox}
